%% BLOCK 3 -- Introduction to Machine Learning Problem Framing.
%% Fragment of ai_foundations.tex; not a standalone document.
%% Source: developers.google.com/machine-learning/problem-framing, all pages
%% (Problem framing, Understand the problem, Framing an ML problem, Implementing
%% a model, Summary).

\actpage{3}{Problem Framing}{Whether ML is the right approach at all, and if it is, how to state the problem in ML terms.}

% ======================================================================
\sectionsubtitle{Analyzing a problem to isolate what has to be solved.}
\section[Problem framing]{Problem Framing}

\begin{frame}{\modulehead{Problem framing}{45 minutes}}
  \begin{withmargin}
    \begin{tdmain}
      Problem framing is the process of analyzing a problem to isolate the
      individual elements that need to be addressed to solve it. It helps
      determine your project's technical feasibility and provides a clear set of
      goals and success criteria. When considering an ML solution, effective
      problem framing can determine whether or not your product
      \alertbf{ultimately succeeds}.
      \vskip0.8em
      At a high level, ML problem framing consists of two distinct steps:
      determining whether ML is the right approach for solving a problem, and
      framing the problem in ML terms.
    \end{tdmain}
    \begin{tdmargin}
      \vskip0.3em
      {\color{tdfg}\textit{Formal problem framing is the critical beginning for
      solving an ML problem, as it forces us to better understand both the
      problem and the data in order to design and build a bridge between them.}}
      \par\smallskip
      \hfill{}TensorFlow engineer
      \par\medskip
      Problem framing ensures that an ML approach is a good solution to the
      problem \gterm{before} beginning to work with data and train a model.
    \end{tdmargin}
  \end{withmargin}
  \slidesources{GoogleProblemFraming2025}
\end{frame}

\begin{frame}{State the goal, in non-ML terms}
  \begin{withmargin}
    \begin{tdmain}
      Begin by stating your goal in non-ML terms. The goal is the answer to the
      question, \gterm{what am I trying to accomplish?}
      \vskip0.5em
      \footnotesize
      \begin{tabular}{@{}p{0.26\linewidth}p{0.66\linewidth}@{}}
        \toprule
        \textbf{Application} & \textbf{Goal}\\
        \midrule
        Weather app & Calculate precipitation in six-hour increments for a geographic region\\
        Fashion app & Generate a variety of shirt designs\\
        Video app & Recommend useful videos\\
        Mail app & Detect spam\\
        Financial app & Summarize financial information from multiple news sources\\
        Map app & Calculate travel time\\
        Banking app & Identify fraudulent transactions\\
        Dining app & Identify cuisine by a restaurant's menu\\
        Ecommerce app & Reply to reviews with helpful answers\\
        \bottomrule
      \end{tabular}
    \end{tdmain}
    \begin{tdmargin}
      {\color{tdfg}\textbf{Objectives.}}
      \par\medskip
      Identify if ML is a good solution for a problem. Learn how to frame an ML
      problem. Understand how to pick the right model and define success
      metrics.
    \end{tdmargin}
  \end{withmargin}
\end{frame}

\begin{frame}{ML is a specialized tool}
  \begin{withmargin}
    \begin{tdmain}
      Some view ML as a universal tool that can be applied to all problems. In
      reality, ML is a \alertbf{specialized tool suitable only for particular
      problems}. You don't want to implement a complex ML solution when a
      simpler non-ML solution will work.
      \vskip0.6em
      \footnotesize
      \begin{tabular}{@{}p{0.20\linewidth}p{0.34\linewidth}p{0.38\linewidth}@{}}
        \toprule
         & \textbf{Output} & \textbf{Training technique}\\
        \midrule
        Predictive ML & Makes a prediction, for example classifying an email as
          spam, guessing tomorrow's rainfall, or predicting a stock price. The
          output can typically be verified against reality.
        & Typically uses lots of data to train a supervised, unsupervised, or
          reinforcement learning model to perform a specific task.\\[0.4em]
        Generative AI & Generates output based on the user's intent, for example
          summarizing an article or producing an audio clip.
        & Typically uses lots of unlabeled data to train a large language model
          or image generator to fill in missing data.\\
        \bottomrule
      \end{tabular}
    \end{tdmain}
    \begin{tdmargin}
      Both take text, image, audio, video, or numerical input.
      \par\medskip
      A generative model can then be used for tasks that can be framed as
      fill-in-the-blank tasks, or fine-tuned by training on labeled data for a
      specific task, like classification.
    \end{tdmargin}
  \end{withmargin}
\end{frame}

\begin{frame}{Benchmark against a non-ML solution}
  \begin{withmargin}
    \begin{tdmain}
      To confirm that ML is the right approach, first verify that your current
      non-ML solution is optimized. If you don't have one, try solving the
      problem manually using a heuristic. The non-ML solution is the
      \gterm{benchmark} you'll use to determine whether ML is a good use case.
      \vskip0.7em
      \textbf{Quality.} How much better do you think an ML solution can be? If
      you think it might be only a small improvement, that might indicate the
      current solution is the best one.
      \vskip0.5em
      \textbf{Cost and maintenance.} How expensive is the ML solution in both
      the short and long term?
    \end{tdmain}
    \begin{tdmargin}
      Can the ML solution justify the increase in cost? Small improvements in
      large systems can easily justify the cost.
      \par\medskip
      How much maintenance will the solution require? In many cases, ML
      implementations need dedicated long-term maintenance.
      \par\medskip
      Does your product have the resources to support training or hiring people
      with ML expertise?
    \end{tdmargin}
  \end{withmargin}
\end{frame}

\begin{frame}{Six characteristics of usable data}
  \begin{withmargin}
    \begin{tdmain}
      \footnotesize
      \begin{itemize}
        \item \textbf{Abundant.} The more relevant and useful examples in your
          dataset, the better your model will be.
        \item \textbf{Consistent and reliable.} An ML-based weather model
          benefits from data gathered over many years from the same reliable
          instruments.
        \item \textbf{Trusted.} Will the data come from trusted sources you
          control, like logs from your product, or from sources you have little
          insight into, like the output of another ML system?
        \item \textbf{Available.} Make sure all inputs are available at
          prediction time in the correct format. If a feature value will be
          difficult to obtain then, \gterm{omit that feature}.
        \item \textbf{Correct.} Some labels will inevitably be wrong, but if
          more than a small percentage are, the model will produce poor
          predictions.
        \item \textbf{Representative.} The datasets should accurately reflect
          the events, user behaviors, and phenomena of the real world.
      \end{itemize}
    \end{tdmain}
    \begin{tdmargin}
      \vfill
      Data is the driving force of predictive ML. If you can't get the data you
      need in the required format, your model will make poor predictions.
      \vfill
    \end{tdmargin}
  \end{withmargin}
\end{frame}

\begin{frame}{Predictive power}
  \begin{withmargin}
    \begin{tdmain}
      The more correlated a feature is with a label, the more likely it is to
      predict it. In a weather dataset, \texttt{cloud\_coverage},
      \texttt{temperature}, and \texttt{dew\_point} would be better predictors
      of rain than \gterm{\texttt{moon\_phase}} or
      \gterm{\texttt{day\_of\_week}}.
      \vskip0.8em
      Determining which features have predictive power can be a time consuming
      process. You can explore it manually by removing and adding a feature
      while training. You can automate it using algorithms such as Pearson
      correlation, adjusted mutual information, and Shapley value, which provide
      a numerical assessment.
    \end{tdmain}
    \begin{tdmargin}
      For the video app, you could hypothesize that \texttt{video\_description},
      \texttt{length}, and \texttt{views} might be good predictors of which
      videos a user would want to watch.
      \par\medskip
      Three key attributes to look for when analyzing your datasets:
      representative of the real world, contains correct values, and features
      have predictive power for the label.
    \end{tdmargin}
  \end{withmargin}
\end{frame}

\begin{frame}{Predictions versus actions}
  \vfill
  \begin{takeaway}
    There's no value in predicting something\\
    if you can't turn the prediction into an \alertbf{action}\\
    that helps users.
  \end{takeaway}
  \vfill
  {\small A model that predicts whether a user will find a video useful should
  feed into an app that recommends useful videos. A model that predicts whether
  it will rain should feed into a weather app.}
  \keyterms{dataset; feature; heuristic; label; prediction}
\end{frame}

% ======================================================================
\sectionsubtitle{Ideal outcome, model output, success metrics.}
\section[ML framing]{Framing an ML Problem}

\begin{frame}{Ideal outcome and the model's goal}
  \begin{withmargin}
    \begin{tdmain}
      \footnotesize
      \begin{tabular}{@{}p{0.16\linewidth}p{0.36\linewidth}p{0.40\linewidth}@{}}
        \toprule
        \textbf{App} & \textbf{Ideal outcome} & \textbf{Model's goal}\\
        \midrule
        Weather & Calculate precipitation in six hour increments for a region
          & Predict six-hour precipitation amounts for specific regions\\[0.3em]
        Fashion & Generate a variety of shirt designs
          & Generate three varieties of a shirt design from text and an image\\[0.3em]
        Video & Recommend useful videos & Predict whether a user will click on a video\\[0.3em]
        Mail & Detect spam & Predict whether or not an email is spam\\[0.3em]
        Map & Calculate travel time & Predict how long it will take to travel between two points\\[0.3em]
        Banking & Identify fraudulent transactions & Predict if a transaction was made by the card holder\\[0.3em]
        Ecommerce & Generate customer support replies & Generate replies using sentiment analysis and the knowledge base\\
        \bottomrule
      \end{tabular}
    \end{tdmain}
    \begin{tdmargin}
      Independent of the ML model, what's the ideal outcome? What is the exact
      task you want your product or feature to perform? This is the same
      statement you defined when you stated the goal.
      \par\medskip
      Then \gterm{connect} the model's goal to the ideal outcome by explicitly
      defining what you want the model to do.
    \end{tdmargin}
  \end{withmargin}
\end{frame}

\begin{frame}{Identify the output you need}
  \begin{withmargin}
    \begin{tdmain}
      If you need to classify something or make a numeric prediction, you'll
      probably use \textbf{predictive ML}. If you need to generate new content
      or produce output related to natural language understanding, you'll
      probably use \textbf{generative AI}.
      \vskip0.5em
      \footnotesize
      \begin{tabular}{@{}p{0.26\linewidth}p{0.32\linewidth}p{0.34\linewidth}@{}}
        \toprule
        \textbf{ML system} & & \textbf{Example output}\\
        \midrule
        Classification & Binary & Classify an email as spam or not spam\\
         & Multiclass single-label & Classify an animal in an image\\
         & Multiclass multi-label & Classify all the animals in an image\\[0.3em]
        Numerical & Unidimensional regression & Predict the number of views a video will get\\
         & Multidimensional regression & Predict blood pressure, heart rate, and cholesterol\\
        \bottomrule
      \end{tabular}
    \end{tdmain}
    \begin{tdmargin}
      \gterm{Generative AI outputs.} Text: summarize an article, reply to
      reviews, translate documents, write product descriptions, analyze legal
      documents. Image: marketing images, visual effects, product design
      variations. Audio: dialogue in a specific accent, a short musical
      composition. Video: realistic-looking videos. Multimodal: a video with
      text captions.
    \end{tdmargin}
  \end{withmargin}
\end{frame}

\begin{frame}{Classification or regression: the caching example}
  \begin{withmargin}
    \begin{tdmain}
      You want to cache videos based on predicted popularity: 50 or more views
      uses the expensive cache, 30 to 50 the cheap cache, fewer than 30 no
      cache. A regression model seems right, since you're predicting a number.
      \vskip0.8em
      But when training it, you realize it produces the \gterm{same loss} for a
      prediction of 28 and 32 on a video that got 30 views. Your app will behave
      very differently for 28 versus 32, but the model considers both
      predictions equally good.
      \vskip0.5em
      \footnotesize
      A classification model would produce a higher loss for 28 than 32. In a
      sense, classification models produce thresholds by default.
    \end{tdmain}
    \begin{tdmargin}
      \textbf{Regression models are unaware of product-defined thresholds.} If
      your app's behavior changes significantly because of small differences in
      a regression model's predictions, consider a classification model instead.
    \end{tdmargin}
  \end{withmargin}
\end{frame}

\begin{frame}{Two lessons from the caching example}
  \begin{withmargin}
    \begin{tdmain}
      \textbf{Predict the decision.} When possible, predict the decision your
      app will take. In the video example, a classification model would predict
      the decision if the categories it classified videos into were "no cache",
      "cheap cache", and "expensive cache". \gterm{Hiding your app's behavior
      from the model can cause your app to produce the wrong behavior.}
      \vskip0.7em
      \textbf{Understand the problem's constraints.} If your app takes different
      actions based on different thresholds, determine if those thresholds are
      fixed or dynamic.
    \end{tdmain}
    \begin{tdmargin}
      \textbf{Dynamic thresholds:} use a regression model and set the threshold
      limits in your app's code. This lets you easily update the thresholds
      while still having the model make reasonable predictions.
      \par\medskip
      \textbf{Fixed thresholds:} use a classification model and label your
      datasets based on the threshold limits.
      \par\medskip
      Most cache provisioning is dynamic, so for that problem regression is the
      best choice. For many problems the thresholds are fixed, making
      classification best.
    \end{tdmargin}
  \end{withmargin}
\end{frame}

\begin{frame}{Proxy labels}
  \begin{withmargin}
    \begin{tdmain}
      In the weather app the label directly addresses the ideal outcome:
      \texttt{precipitation\_amount}. In the video app the ideal outcome is to
      recommend useful videos, but there's no label called
      \texttt{useful\_to\_user}. Proxy labels substitute for labels that aren't
      in the dataset.
      \vskip0.5em
      \footnotesize
      \begin{tabular}{@{}p{0.42\linewidth}p{0.50\linewidth}@{}}
        \toprule
        \textbf{Proxy label} & \textbf{Issue}\\
        \midrule
        User will click "like" & Most users never click "like"\\
        Video will be popular & Not personalized\\
        User will share the video & Some users don't share; sometimes people share videos because they dislike them\\
        User will click play & Maximizes clickbait\\
        How long they watch & Favors long videos\\
        How many times they rewatch & Favors rewatchable genres\\
        \bottomrule
      \end{tabular}
    \end{tdmain}
    \begin{tdmargin}
      A workable proxy label for usefulness might be
      \texttt{shared OR liked}.
      \par\medskip
      \gterm{No proxy label can be a perfect substitute for your ideal outcome.
      All will have potential problems. Pick the one that has the least problems
      for your use case.}
    \end{tdmargin}
  \end{withmargin}
\end{frame}

\begin{frame}[fragile]{Generation: three ways to customize}
  \begin{withmargin}
    \begin{tdmain}
      In most cases you won't train your own generative model, because that
      requires massive amounts of training data and computational resources.
      Instead you'll customize a pre-trained one.
      \vskip0.5em
      \footnotesize
      \textbf{Distillation} creates a smaller version of a larger model: you
      generate a synthetic labeled dataset from the larger model and use it to
      train the smaller one.
      \vskip0.3em
      \textbf{Fine-tuning or parameter-efficient tuning} further trains the
      model on a dataset containing examples of the type of output you want.
      \vskip0.3em
      \textbf{Prompt engineering} tells the model the task you want it to do, or
      shows illustrative examples with the desired outputs.
    \end{tdmain}
    \begin{tdmargin}
      \begin{tdcode}
\begin{Verbatim}[fontsize=\tiny]
Translate words from English
to Spanish.

English: Car
Spanish: Auto

English: Airplane
Spanish: Avion

English: Home
Spanish: ______
\end{Verbatim}
      \end{tdcode}
      \par\smallskip
      Distillation and fine-tuning update the model's parameters. Prompt
      engineering \gterm{does not}.
    \end{tdmargin}
  \end{withmargin}
\end{frame}

\begin{frame}{Success metrics}
  \begin{withmargin}
    \begin{tdmain}
      Success metrics define what you care about, like engagement or helping
      users take appropriate action. They \alertbf{differ from the model's
      evaluation metrics}, like accuracy, precision, recall, or AUC.
      \vskip0.5em
      \footnotesize
      \begin{tabular}{@{}p{0.14\linewidth}p{0.78\linewidth}@{}}
        \toprule
        \multicolumn{2}{@{}l}{\textbf{Weather app}}\\
        Success & Users open the "Will it rain?" feature 50 percent more often than before\\
        Failure & Users open the feature no more often than before\\
        \midrule
        \multicolumn{2}{@{}l}{\textbf{Video app}}\\
        Success & Users spend on average 20 percent more time on the site\\
        Failure & Users spend on average no more time on site than before\\
        \bottomrule
      \end{tabular}
      \vskip0.4em
      Define ambitious success metrics. High ambitions can cause gaps between
      success and failure. The undefined gap is not what's important.
    \end{tdmain}
    \begin{tdmargin}
      What's important is your model's capacity to move closer to, or exceed,
      the definition of success.
      \par\medskip
      \gterm{Would improving the model get you closer to your defined success
      criteria?} A model might have great evaluation metrics but not move you
      closer, indicating that even a perfect model wouldn't meet your criteria.
    \end{tdmargin}
  \end{withmargin}
\end{frame}

\begin{frame}{Is the model worth improving?}
  \begin{withmargin}
    \begin{tdmain}
      \begin{itemize}
        \item \textbf{Not good enough, but continue.} The model shouldn't be
          used in production, but over time it might be significantly improved.
        \item \textbf{Good enough, and continue.} The model could be used in
          production, and it might be further improved.
        \item \textbf{Good enough, but can't be made better.} The model is in
          production, but it is probably as good as it can be.
        \item \textbf{Not good enough, and never will be.} The model shouldn't
          be used in production and no amount of training will probably get it
          there.
      \end{itemize}
      \vskip0.5em
      \footnotesize
      When deciding to improve the model, re-evaluate if the increase in
      resources, like engineering time and compute costs, \gterm{justify} the
      predicted improvement.
    \end{tdmain}
    \begin{tdmargin}
      After defining the success and failure metrics, determine how often you'll
      measure them: six days, six weeks, or six months after implementing the
      system.
      \par\medskip
      When analyzing failure metrics, determine \textbf{why} the system failed.
      The model might be predicting which videos users will click, but start
      recommending clickbait titles that cause engagement to drop off.
    \end{tdmargin}
  \end{withmargin}
  \keyterms{accuracy; AUC; classification model; precision; proxy label;
    regression model}
\end{frame}

% ======================================================================
\sectionsubtitle{Start simple, then decide whether complexity is justified.}
\section[Implementing]{Implementing a Model}

\begin{frame}{Start simple}
  \begin{withmargin}
    \begin{tdmain}
      Most of the work in ML is on the \gterm{data side}, so getting a full
      pipeline running for a complex model is harder than iterating on the model
      itself. After setting up your data pipeline and implementing a simple
      model that uses a few features, you can iterate on creating a better
      model.
      \vskip0.8em
      Simple models provide a good baseline, even if you don't end up launching
      them. In fact, using a simple model is probably better than you think.
      Starting simple helps you determine whether or not a complex model is
      \alertbf{even justified}.
    \end{tdmain}
    \begin{tdmargin}
      Trained models only really work when the label and features match your
      dataset exactly. If a trained model uses 25 features and your dataset
      includes only 24 of them, it will most likely make bad predictions.
      \par\medskip
      Commonly, practitioners use matching subsections of inputs from a trained
      model for fine-tuning or transfer learning.
      \par\medskip
      If your solution is a generative AI model, you'll almost always fine-tune
      a pre-trained model instead of training your own.
    \end{tdmargin}
  \end{withmargin}
\end{frame}

\begin{frame}{Monitoring, decided during framing}
  \begin{withmargin}
    \begin{tdmain}
      \textbf{Model deployment.} A newly trained model might be worse than the
      model currently in production. If it is, you'll want to prevent it from
      being released and get an alert that your automated deployment has failed.
      \vskip0.6em
      \textbf{Training-serving skew.} If any incoming features used for
      inference have values outside the distribution range of the training data,
      you'll want to be alerted, because the model will likely make poor
      predictions.
      \vskip0.6em
      \textbf{Inference server.} If you're providing inferences through an RPC
      system, monitor the RPC server itself and get an alert if it stops
      providing inferences.
    \end{tdmain}
    \begin{tdmargin}
      If your model was trained to predict temperatures for equatorial cities at
      sea level, your serving system should alert you of incoming data with
      latitudes, longitudes, or altitudes \gterm{outside the range} the model
      was trained on.
      \par\medskip
      Conversely, it should alert you if the model is making predictions outside
      the distribution range seen during training.
    \end{tdmargin}
  \end{withmargin}
\end{frame}

\begin{frame}{The two steps, in six moves}
  \begin{withmargin}
    \begin{tdmain}
      \textbf{Verify that ML is a good approach.}
      \vskip0.3em
      \begin{itemize}
        \item Understand the problem.
        \item Identify a clear use case.
        \item Understand the data.
      \end{itemize}
      \vskip0.7em
      \textbf{Frame the problem in ML terms.}
      \vskip0.3em
      \begin{itemize}
        \item Define the ideal outcome and the model's goal.
        \item Identify the model's output.
        \item Define success metrics.
      \end{itemize}
    \end{tdmain}
    \begin{tdmargin}
      These steps save time and resources by setting \gterm{clear goals} and
      providing a shared framework for working with other ML practitioners.
      \par\medskip
      When implementing ML solutions, always follow Google's Responsible AI
      Principles. For a hands-on introduction to improving fairness and
      mitigating bias, see the ML fairness module in block 2.
    \end{tdmargin}
  \end{withmargin}
  \slidesources{GoogleProblemFraming2025}
\end{frame}
